\section{Energy Consumption}

Energy consumption patterns in Indonesia reflect the country's unique economic structure and development trajectory. As a nation comprising over 17,000 islands with a population exceeding 270 million, Indonesia faces distinct challenges in energy provision and management. The industrial sector accounts for approximately 44.2\% of total final energy consumption, followed by the transport sector at 40.3\% \cite{IEA2021}. This distribution highlights the energy-intensive nature of Indonesia's manufacturing and logistics operations, which are crucial to its export-oriented economy.

The electricity generation mix remains heavily reliant on fossil fuels, with coal contributing 60.3\% of the total generation capacity in 2020 \cite{PLN2021}. This dependence on coal has positioned Indonesia as one of the world's largest coal consumers, with annual consumption reaching 126.2 million tonnes of oil equivalent (Mtoe) \cite{BP2021}. The government's commitment to achieving a 23\% renewable energy share by 2025 presents significant challenges, as current renewable energy contribution stands at only 11.2\% \cite{MEMR2021}.

Industrial energy consumption patterns vary significantly across different sub-sectors. The manufacturing industry, particularly energy-intensive sectors such as cement, steel, and pulp and paper, accounts for a substantial portion of industrial energy use. The cement industry alone consumes approximately 5.5 Mtoe annually, with energy costs representing 30-40\% of production costs \cite{CementAssociation2020}. Many industrial facilities have begun implementing energy efficiency measures, with the potential for energy savings estimated at 10-30\% through technological improvements and operational optimization \cite{ESDM2019}.

The commercial and residential sectors contribute approximately 15\% and 14\% respectively to total final energy consumption \cite{IEA2021}. Rapid urbanization and rising living standards have driven increased electricity demand in these sectors, with residential electricity consumption growing at an average annual rate of 6.2\% between 2010 and 2020 \cite{PLN2021}. The widespread use of air conditioning, particularly in urban areas, has become a significant driver of peak electricity demand.

Transportation energy consumption presents unique challenges in the Indonesian context. The dominance of road transport, which accounts for 85\% of passenger transport and 90\% of freight transport \cite{BPS2021}, has led to increasing fuel consumption. The transport sector consumed 71.4 Mtoe of energy in 2020, with road transport accounting for 95\% of this total \cite{IEA2021}. The government's fuel subsidy policy, while aimed at maintaining affordability, has created market distortions and reduced incentives for energy efficiency improvements.

Energy intensity, measured as energy consumption per unit of GDP, has shown a declining trend in recent years, falling from 0.58 kgoe/2010 USD in 2010 to 0.45 kgoe/2010 USD in 2020 \cite{IEA2021}. This improvement reflects both structural changes in the economy and efficiency gains in various sectors. However, the energy intensity of Indonesia's industrial sector remains higher than the global average, indicating significant potential for improvement through technology adoption and process optimization.

The implementation of energy management systems has gained momentum, particularly among large industrial consumers. The adoption of ISO 50001 energy management standards has increased, with over 50 certified facilities across various sectors \cite{BSI2021}. These systems have demonstrated energy savings of 10-20\% in certified facilities, contributing to both cost reduction and emission mitigation \cite{ESDM2019}.

Renewable energy development faces several challenges in Indonesia, including geographic constraints, infrastructure limitations, and financing barriers. The country's vast archipelago nature makes grid connection and energy distribution challenging, particularly for remote areas. Despite these challenges, significant progress has been made in geothermal energy development, with Indonesia holding 40\% of the world's geothermal reserves \cite{MEMR2021}. Current installed geothermal capacity stands at 2.3 GW, with plans to expand to 7.2 GW by 2025 \cite{PLN2021}.

Energy efficiency policies and programs have been implemented at both national and local levels. The National Energy Conservation Master Plan (RIKEN) aims to reduce energy intensity by 1\% annually through various measures, including building energy codes, industrial energy efficiency programs, and public awareness campaigns \cite{ESDM2020}. The success of these programs varies across regions, with more developed areas showing higher adoption rates of energy-efficient technologies.

The COVID-19 pandemic has had a significant impact on energy consumption patterns, with industrial and commercial energy use declining by 8-12\% in 2020 \cite{IEA2021}. However, the residential sector saw increased energy consumption due to work-from-home arrangements and increased use of home appliances. The pandemic has also accelerated the adoption of digital technologies for energy management and monitoring, particularly in industrial facilities.

Looking forward, Indonesia's energy consumption is projected to grow at an average annual rate of 4.2\% until 2030 \cite{MEMR2021}. This growth will be driven by continued economic development, urbanization, and rising living standards. The successful implementation of energy efficiency measures and renewable energy development will be crucial in managing this growth sustainably. The government's commitment to achieving net-zero emissions by 2060 will require significant transformations in energy consumption patterns across all sectors.

The industrial sector's role in energy transition is particularly critical, given its substantial energy consumption and emission footprint. Many large companies have begun setting science-based targets for emission reduction, with energy efficiency and renewable energy procurement being key strategies. The development of industrial symbiosis and circular economy approaches offers additional opportunities for energy optimization and waste reduction.

Energy consumption patterns in Indonesia reflect the complex interplay between economic development, technological capabilities, and policy frameworks. While significant progress has been made in improving energy efficiency and developing renewable energy, substantial challenges remain in achieving sustainable energy consumption patterns. The successful management of energy consumption will be crucial for Indonesia's economic development and environmental sustainability goals.